%%%%%%%%%%%%%%%%%%%%%%%%%%%%%%%%%%%%%%%%%
% Structured General Purpose Assignment
% LaTeX Template
%
% This template has been downloaded from:
% http://www.latextemplates.com
%
% Original author:
% Ted Pavlic (http://www.tedpavlic.com)
%
%%%%%%%%%%%%%%%%%%%%%%%%%%%%%%%%%%%%%%%%%

%----------------------------------------------------------------------------------------
%   PACKAGES AND OTHER DOCUMENT CONFIGURATIONS
%----------------------------------------------------------------------------------------

\documentclass{article}

\usepackage{fancyhdr} % Required for custom headers
\usepackage{lastpage} % Required to determine the last page for the footer
\usepackage{extramarks} % Required for headers and footers
\usepackage{graphicx} % Required to insert images
\usepackage{lipsum} % Used for inserting dummy 'Lorem ipsum' text into the template

\usepackage[T1]{fontenc} % Codificación de las fuentes utilizadas
\usepackage[spanish]{babel} % Español como idioma principal del texto (permite hyphenation de palabras al final de una línea)
\selectlanguage{spanish}
\usepackage{hyperref}
\usepackage{listings}

\usepackage[T1]{fontenc} % Codificación de las fuentes utilizadas
\usepackage[spanish]{babel} % Español como idioma principal del texto (permite hyphenation de palabras al final de una línea)


\usepackage{graphicx}
\usepackage{url}

\graphicspath{{Figures/}{Diagrams}{Chapters/}}  % Location of the graphics files (set up for graphics to be in PDF format)

\selectlanguage{spanish}

\setcounter{tocdepth}{1}

% Include any extra LaTeX packages required
\usepackage[square, numbers, comma, sort&compress]{natbib}  % Use the "Natbib" style for the references in the Bibliography
\usepackage{verbatim}  % Needed for the "comment" environment to make LaTeX comments
\usepackage{vector}  % Allows "\bvec{}" and "\buvec{}" for "blackboard" style bold vectors in maths
\hypersetup{urlcolor=blue, colorlinks=true}  % Colours hyperlinks in blue, but this can be distracting if there are many links.
\usepackage{hyperref}
% \usepackage[pdfauthor={Diego Martín Arroyo},
%             pdftitle={Diseño e implementación de un sistema de computación distribuida con
% Raspberry Pi, y estudio comparativo del mismo frente a otras soluciones},
%             pdfsubject={Memora del Trabajo de Fin de Grado},
%             pdfproducer={XeLaTeX with hyperref},
%             pdfcreator={XeLaTeX},
%             pdfkeywords={Computación Paralela, Sistema Distribuido, Raspberry}
%             ]{hyperref}
%% ----------------------------------------------------------------

%% --------------------------------------------------------------------------------------------------------------------------------
%http://tex.stackexchange.com/a/85218/76599
\usepackage{fancyvrb}
\usepackage[dvipsnames]{xcolor}

% redefine \VerbatimInput
\RecustomVerbatimCommand{\VerbatimInput}{VerbatimInput}% Inclusión de archivos de texto plano
{fontsize=\footnotesize,
 %
 frame=lines,  % top and bottom rule only
 framesep=2em, % separation between frame and text
 rulecolor=\color{Gray},
 %
 label=\fbox{\color{Black}data.txt},
 labelposition=topline,
 %
 commandchars=\|\(\), % escape character and argument delimiters for
                      % commands within the verbatim
 commentchar=*        % comment character
}

\usepackage{listings} % Requerido para la inserción de código
%Listings command

\usepackage{float}
\newcommand*\lstinputpath[1]{\lstset{inputpath=#1}}
\lstinputpath{Code/}

\newcounter{undefinedreferences}
\setcounter{undefinedreferences}{0}

\newcommand{\citationneeded}[1][None]{\stepcounter{undefinedreferences}\textsuperscript{\color{blue} [Citation needed: #1]}}

\newcommand{\checkreferences}{
\ifnum\value{undefinedreferences} > 0
\begin{center}
\immediate\write18{wget -O Figures/protester.png -nc http://imgs.xkcd.com/comics/wikipedian_protester.png}
\includegraphics[width=\textwidth]{protester.png}\\
There are \arabic{undefinedreferences} undefined references
\end{center}
\else
No undefined references. Good!
\fi
}


%https://github.com/pads-fhs/LaTeX-Template-Thesis/blob/master/lststyles.tex
\lstdefinelanguage{JavaScript}{
  keywords={typeof, new, true, false, catch,%
    function, return, null, catch, switch, var,%
    if, in, while, do, else, case, break},
  ndkeywords={class, export, boolean, throw, implements, import, this},
  sensitive=false,
  comment=[l]{//},
  morecomment=[s]{/*}{*/},
  morestring=[b]',
  morestring=[b]"
}
\newcommand{\lstsetjavascript}{
  \lstset{
		language=JavaScript,
		breaklines=true,
		commentstyle=\textit,
		basicstyle=\ttfamily,
		keywordstyle=\bfseries,
		stringstyle=\ttfamily,
		showstringspaces=false,
		frame=single,
		tabsize=2
  }
}

\lstdefinelanguage{log}{
  keywords={typeof, new, true, false, catch,%
    function, return, null, catch, switch, var,%
    if, in, while, do, else, case, break},
  ndkeywords={class, export, boolean, throw, implements, import, this},
  sensitive=false,
  comment=[l]{//},
  morecomment=[s]{/*}{*/},
  morestring=[b]',
  morestring=[b]"
}
\newcommand{\lstsetlog}{
  \lstset{
		language=log,
		breaklines=true,
		commentstyle=\textit,
		basicstyle=\ttfamily,
		keywordstyle=\bfseries,
		stringstyle=\ttfamily,
		showstringspaces=false,
		frame=single,
		tabsize=2
  }
}

\lstloadlanguages{Java,XML, JavaScript, log}

\newcommand{\javascriptcode}[4]{
	\lstinputlisting[caption=#2,label=#1, firstline=#3, lastline=#4]{#1.json}
}

\newcommand{\logcode}[4]{
	\lstinputlisting[caption=#2,label=#1, firstline=#3, lastline=#4]{#1.log}
}

\usepackage[bottom]{footmisc} %The footnotes go at the bottom of t\usepackage{dtklogos}he page, instead next to the last line.
%Ajustes para Java
% \lstset{
% 	language=java,
%  	frame=single, % Un marco simple alrededor del código
%     basicstyle=\small\ttfamily, % Utilizar fuente true type pequeña
%     keywordstyle=[1]\color{Blue}\bf, % Funciones en negrita y azul
%     keywordstyle=[2]\color{Purple}, % Argumentos en morado
%     keywordstyle=[3]\color{Blue}\underbar, % Funciones personalizadas subrayadas en azul
%     identifierstyle=, % Nada especial acerca de identificadores
%     commentstyle=\usefont{T1}{pcr}{m}{sl}\color{Green}\small, % Los comentarios se renderizan en fuente pequeña verde
%     stringstyle=\color{Purple}, % Cadenas en morado
%     showstringspaces=false, % No se muestran los espacios entre cadenas
%     tabsize=5, % 5 espacios por tabulado
%     %
%     % Put standard Perl functions not included in the default language here
%     %morekeywords={rand},
%     %
%     % Put Perl function parameters here
%     %morekeywords=[2]{on, off, interp},
%     %
%     % Put user defined functions here
%     %morekeywords=[3]{test},\usepackage{dtklogos}
%    	%
%     morecomment=[l][\color{Blue}]{...}, % Line continuation (...) like blue comment
%     numbers=left, % Número de línea a la izquierda
%     firstnumber=1, % Número de línea comienza en 1
%     numberstyle=\tiny\color{Blue}, % Los números de línea son azules y pequeños
%     stepnumber=5, % Los números de línea van de 5 en 5
%     breaklines=true % Salto de línea si el texto no entra. See http://stackoverflow.com/a/1875803
% }

%\usepackage{xltxtra} % XeLaTeX logo. Yep, just that
%http://tex.stackexchange.com/a/73179/76599
\usepackage{metalogo}
%\usepackage{dtklogos} %BibTeX logo
\def\BibTeX{{\rm B\kern-.05em{\sc i\kern-.025em b}\kern-.08em
    T\kern-.1667em\lower.7ex\hbox{E}\kern-.125emX}}

\newenvironment{alignedDescription}[2][0pt]
  {\begin{list}{}%
    {\renewcommand\makelabel[1]{\textsf{\textbf{##1}}\hfil}%
     \settowidth\labelwidth{\makelabel{#2}}%
     \setlength\leftmargin{\labelwidth+\labelsep + #1}}}%
  {\end{list}}

\newenvironment{elements}
{\begin{quote}\itshape\centering\small}
{\end{quote}}

\newenvironment{cabstract}
{\begin{quote}\itshape\centering\small}
{\end{quote}}

%\usepackage[xindy]{glossaries}
%\newcommand{\chapterabstract}{1}{
%	\begin{center}
%	\small\textit
%	#1
%	\end{center}
%}

\usepackage{xcolor,colortbl}
\usepackage{float}
%----------------------------------------------------------------------------------------
%   NAME AND CLASS SECTION
%----------------------------------------------------------------------------------------

\newcommand{\hmwkTitle}{Evaluación del rendimiento del sistema en un entorno real} % Assignment title
\newcommand{\hmwkDueDate}{Domingo,\ 25\ de\ mayo\ de\ 2015} % Due date
\newcommand{\hmwkAuthorName}{Diego Martín Arroyo} % Your name

\graphicspath{{Figures/}}
%----------------------------------------------------------------------------------------
%   TITLE PAGE
%----------------------------------------------------------------------------------------

\title{\hmwkTitle}
\author{\textbf{\hmwkAuthorName}}
\date{\hmwkDueDate}

%----------------------------------------------------------------------------------------

\begin{document}

\maketitle
\begin{abstract}
Una adecuada evaluación de las herramientas, protocolos y arquitectura planteadas debe incluir un conjunto de pruebas en entornos en los que la versión final del sistema debe integrarse. Dichas pruebas deben ser preparadas de antemano.
\end{abstract}

%----------------------------------------------------------------------------------------
%   TABLE OF CONTENTS
%----------------------------------------------------------------------------------------

\setcounter{tocdepth}{1}

\tableofcontents
\newpage

\section{Marcobootstrap}

Se realizarán las siguientes pruebas

\subsection*{Instalación del sistema operativo en un nodo sin asistencia del administrador, evaluando el tiempo de descarga e instalación}

Sistemas utilizados
\begin{itemize}
\item Una Rasbperry Pi Modelo 2B
\item Un servidor Sun
\item Imagen creada con las versiones estables más recientes de todas las herramientas, con un tamaño de 
\item La prueba se realiza en la infraestructura final, conectando los nodos a las tomas de red del Laboratorio de Ingeniería Informática.
\end{itemize}

Resultados
\begin{itemize}
\item \textbf{Tiempo de descarga}: 3 minutos, 27 segundos
\item \textbf{Tiempo de instalación}: 11 minutos, 49 segundos
\item \textbf{Tiempo de instalación}: 22 minutos, 37 segundos
\end{itemize}

\subsection*{Actualización del sistema operativo}

Se realiza la instalación de una versión alternativa del sistema operativo, comprobando que los componentes no actualizables permanecen intactas tras el proceso (ver documentación).

\subsection*{Operaciones de administración}

\subsubsection*{Reinicio}

\subsubsection*{Cancelación}

\section{Deployer}

Sistemas utilizados
\begin{itemize}
\item 4 Rasbperry Pis con la última versión de las herramientas a utilizar.
\item La prueba se realiza en la infraestructura final, conectando los nodos a las tomas de red del Laboratorio de Ingeniería Informática.
\end{itemize}

\section{Logger}

\begin{itemize}
\item 4 Raspberry Pis ejecutando un programa que imprime un número elevado de mensajes en un orden dado
\item La prueba se realiza en la infraestructura final, conectando los nodos a las tomas de red del Laboratorio de Ingeniería Informática.
\end{itemize}

\section{Status monitor}

Se evalúa el rendimiento de la herramienta con los cuatro equipos, uno de ellos compilando el software MPI simultáneamente. El rendimiento de la aplicación es el esperado, contando con una frecuencia de actualización de un segundo de forma constante.

\section{Compilador}

Se realiza como prueba la compilación de la biblioteca openmpi (Versión 1.8.4) con los siguientes parámetros:

\texttt{}
\begin{itemize}
\item \textbf{Tiempo de configuración}: 15 minutos
\item \textbf{Tiempo de compilación}: 23 minutos
\item \textbf{Tiempo de instalación}: 6 minutos
\end{itemize}

\section{Prácticas}

\section{Gestión de usuarios}

\section{LEDs}

\section{MPI}

Se han realizado las siguientes pruebas de rendimiento y viabilidad de MPI:

\subsection{Práctica de la asignatura Arquitectura de Computadores}

Se ha ejecutado la solución a la práctica realizada en la asignatura Arquitectura de Computadores durante el curso 2013-2014.

El enunciado de dicha práctica es el siguiente:

Se desea realizar una transformación de un programa que realiza una serie de cálculos de manera secuencial. En este se quieren ``adivinar" una serie de números de manera aleatoria. Para ello se emplea la siguiente función, a la que se le pasa el número a descubrir y devuelve el número de intentos que ha necesitado para ``adivinarlo" y el tiempo que ha empleado (se ha forzado la semilla a 1 siempre para que el número de intentos para descubrir un mismo número sea siempre el mismo):

SE PIDE: 

Realizar un programa en MPI que permita distribuir el cálculo del ejemplo. Para ello habrá: 
	Un proceso encargado de la E/S que: 

		El proceso de E/S deberá mostrar por pantalla como mínimo: 

			Para cada número a ``adivinar'':
				El número ``adivinado''
				El número de intentos utilizado
				El tiempo empleado

			No es necesario que la información de los números “adivinados” aparezca en el mismo orden que en la lista de números.
		El proceso de E/S mostrará por pantalla el tiempo empleado total y los valores acumulados, tanto de intentos como de tiempos de cálculo. 
		El proceso de E/S puede ser o no calculador. Hay libertad para que el proceso de E/S también calcule o bien sea un mero repartidor de tareas. 

	N procesos calculadores que: 
		Esperarán a que el proceso de E/S le pase un trabajo, ya sea un número o varios números. 

		Para los números “adivinados” deberá pasarle la información necesaria al proceso de E/S para que pueda sacarla por pantalla- 

		Los procesos podrán devolver los resultados como mejor se considere (uno por uno, en una estructura, etc), lo que hay que respetar son los distintos tipos de datos. 

		Al emplear la función rand hay que considerar inicializa la semilla en cada proceso con un valor distinto. El hecho de emplear un procedimiento aleatorio para descubrir los números hará que dependiendo del valor de la semilla los resultados puedan variar. Se ha forzado la semilla a 1 siempre para que el número de intentos para descubrir un mismo número sea siempre el mismo 

	Todos los procesos deberán terminar (no se pueden quedar en un bucle infinito esperando solicitudes). 

Una vez realizado el programa realizar un estudio de rendimiento en el cual se vea que:

	Aumentando el número de procesos se reduce el tiempo de cálculo hasta un cierto momento (por ej.: 1 (secuencial), 2, 4, 8, 16, 64, 128...). 

	Una vez alcanzado el tope de una máquina probar con varias máquinas en red. 

	Es recomendable emplear un número de iteraciones grandes para que los tiempos de cálculo antes de distribuir sean grandes.

\subsubsection{Resultados}

\paragraph{Con 2 nodos}

Intentos Acumulados: 42206944176 
Tiempo Acumulado:    10728.003393   
Tiempo Total:        1691.146348   
Max Intentos:        2798133238 
Max Tiempo:          652.491734

\paragraph{Con 4 nodos}

\section{Backups}

\end{document}
